\section*{Methods}

\subsection*{What makes work avoidable}

This subsection is an existence argument depending on no measurement: a gap must exist between what
a dense-bitmap kernel \emph{spends} and what a cheaper but equally exact algorithm on the same
operands \emph{already achieves}, and representations whose cost is a function of structure rather
than of the universe can enter it. Both compared quantities are achieved; neither lower-bounds the
work the problem demands, and no such bound is claimed here. Every representation is exact and
interconvertible --- \rB, \rS, \rR, \rW\ and the complement view \rC\ encode the same set and every
cell returns the same cardinality --- so nothing is approximated and only the cost changes. Every
quantity in this subsection is a count of stored items or of operations, never a time; constants that convert
these counts to time are empirical and do not enter the derived results. Each result is stated here with its hypotheses; the proof, the worked
numerical checks and the secondary propositions that support but do not themselves state a result are
in Supplementary Note~\ref{note:proofs-avoidance}, for the results up to and including the size bound, and Supplementary Note~\ref{note:proofs-choice}, for the representation-choice optimisation problem that follows it.

Write $A,B \subseteq [0,m)$, $a=|A|$, $b=|B|$, densities $d_A=a/m$, $d_B=b/m$, word width $w$, vector
width $V$ words per operation. Complementation is an involution on density, $d\mapsto 1-d$, so define
the \emph{effective sparsity} $s_X=\min(d_X,1-d_X)\in[0,\tfrac12]$ and $s=\min(s_A,s_B)$.

\paragraph*{What cardinality alone determines.}
Given only $a$ and $b$ --- the cheapest metadata any selector holds --- the answer is confined to
\begin{equation}
\max(0,\,a+b-m) \;\le\; |A \cap B| \;\le\; \min(a,\,b),
\label{eq:interval}
\end{equation}
and every integer in the range is realised (Supplementary Note~\ref{note:proofs-avoidance}). Its width is the quantity the rest
of this subsection turns on:
\begin{equation}
\min(a,b)-\max(0,a+b-m) \;=\; \min\,(\,a,\;b,\;m-a,\;m-b\,) \;=\; m\,\min(s_A,s_B) \;=\; m\,s .
\label{eq:width}
\end{equation}
The four terms are the sizes of the four lists one could enumerate --- $A$, $B$, $A^{c}$, $B^{c}$ ---
so the uncertainty is the size of the smallest, and one near-empty or near-full side collapses it
alone. Above density one half the floor lifts off zero at $(d_A+d_B-1)m$: at $d_A=d_B=0.51$, two per
cent of the universe must overlap.

\paragraph*{The fixed cost, and therefore the headroom.}
A dense bitmap stores $\lceil m/w\rceil$ words whatever it contains, so $\cell{\rB}{\rB}$ performs
$\lceil m/(wV)\rceil$ vector AND-plus-popcount operations per pair, independently of $a$, $b$, of
structure, and of the answer. Against that, suppose each operand $X$ is available as a bitmap and as a
sorted array of whichever of $X$, $X^{c}$ is smaller --- $\Theta(m\,s_X)$ of space, the choice made
from $a$, $b$, $m$ in $O(1)$; no rank index is needed for this result. Then $|A\cap B|$ is answered in
$O(m\,s)$ membership probes (four cases, one per term of Eq.~\eqref{eq:width}; Supplementary Note~\ref{note:proofs-avoidance}).
Counting descriptor items touched, and taking $w \mid m$,
\begin{equation}
\frac{\text{items touched by } \cell{\rB}{\rB}}{\text{items touched by the cheapest list-based cell}}
\;=\; \frac{\lceil m/w\rceil}{m\,s} \;=\; \frac{1}{w\,s},
\label{eq:headroom}
\end{equation}
which exceeds one whenever $s<1/w$ and grows without bound as either operand leaves density one half
--- the headroom, derived rather than observed, and a comparison of two \emph{achieved} costs, not a
lower bound on any algorithm's work. Two of the four cases need a complement view, which is why \rC\
is part of the design rather than an afterthought for the dense tail. Charging measured per-item
costs --- a random probe against a vector word --- turns it into a crossover in $s$ of
\nm{measured crossover in effective sparsity}.

\paragraph*{From descriptor length to operation count.}
Everything else here counts \emph{stored items}, a fact about representation size; the paper's claim
is about work. One identity carries that step. Let $A$ be a bitmap with an $O(1)$ rank index,
$\mathrm{rank}_A(i)=|A\cap[0,i)|$, and $B$ a run container with maximal runs
$\{[u_i,v_i)\}_{i=1}^{r}$. The runs partition $B$, so
\begin{equation}
|A \cap B| \;=\; \sum_{i=1}^{r}\big(\,\mathrm{rank}_A(v_i)-\mathrm{rank}_A(u_i)\,\big),
\label{eq:rank}
\end{equation}
in exactly $2r$ lookups --- independent of $m$, of $|A|$, of $|B|$, and of the run lengths
(Supplementary Note~\ref{note:proofs-avoidance}). A run of any length is priced as a run of length one, because the kernel never
looks inside it; symmetrically $\cell{\rR}{\rR}$ by merge costs $O(r_A+r_B)$ with no index. This is
what makes a run count a cost parameter and not merely a storage parameter, and it is the only bridge
between the two here: it reaches \emph{operation count}, not time, since per-operation constants still
differ across cells.

\paragraph*{What randomness would predict.}
Let each position be set independently with probability $d$, and write $\kappa(\rho)$ for the number
of stored items representation $\rho$ must touch. The bitmap and array entries are exact; the run and
fill-compressed entries are the leading-order forms of exact Bernoulli expectations (Supplementary Note~\ref{note:proofs-avoidance}), written in
$s=\min(d,1-d)$,
\begin{equation}
\mathbb{E}\,\kappa(\rB) = \Big\lceil \tfrac{m}{w} \Big\rceil, \quad
\mathbb{E}\,\kappa(\rS) = m\,d, \quad
\mathbb{E}\,\kappa(\rR) = m\,d(1-d) = m\,s(1-s), \quad
\mathbb{E}\,\kappa(\rW) = \tfrac{m}{w}\big[\,1-(1-s)^{2w}-s^{2w}\,\big],
\label{eq:expected}
\end{equation}
where $\kappa(\rR)$ counts maximal runs and $\kappa(\rW)$ literal words plus fill groups. None of
these expectations is new. The \rW\ row is the word-aligned expected-size model of Wu, Otoo and
Shoshani \cite{wu2006wah}, whose bracket is identical to the one above under their payload width of
$w-1$ bits per word rather than $w$; the symmetry about $d=\tfrac12$ and the resulting
incompressibility at that density are theirs as well. The \rR\ row is the classical expected run
count for a uniformly random arrangement \cite{mood1940runs}, and the interval bound is
the Fr\'echet--Boole inequality. The density--clustering parameterisation is also due to
\cite{wu2006wah}; the design space it spans has since been mapped empirically for Roaring, WAH and
tree-encoded bitmaps by Lang et al. \cite{lang2020teb}, who report that an uncompressed bitmap reads
faster than any of the three compressed formats over $16 \le f \le 128$ and $0.01 \le d < 1$, and
that they expect ``a performance-optimized implementation to dominate an even larger space'' than
their unoptimized baseline did. That is the headroom this paper measures, stated from the compressed
side. They are restated here in one notation because the pairing matrix needs them as
inputs, not because they are results of this paper. With the
complement view making the best of \rS\ and $\rC\!\circ\!\rS$ cost $m\,s$, every entry is constant in
$s$ or strictly increasing on $[0,\tfrac12]$ and symmetric about $d=\tfrac12$: under this null model the
U-shaped cost curve is a theorem --- about the model and the bitmap's flatness, not about any method.
Randomness gives run containers nothing either, since
$\tfrac12 m s\le\mathbb{E}\,\kappa(\rR)\le m\,s+1$, so choosing \emph{among} \{\rS,
$\rC\!\circ\!\rS$, \rR, \rW\} buys at most a constant on structureless data. \rB\ is the exception and
the exception is the point: $\kappa(\rB)$ is $\lceil m/w\rceil$ at every density, so the spread across
the full matrix is Eq.~\eqref{eq:headroom} and unbounded. It is the envelope, not each cell, that
follows $\Theta(\min(m/w,\,m\,s))$.

\paragraph*{The deviation, which is the contribution.}
Real data is not random, and the representations part company with Eq.~\eqref{eq:expected} in a way
density cannot express. At fixed density $d\le\tfrac12$ with $k=dm$, $\kappa(\rB)$ and $\kappa(\rS)$
are constants, while $\kappa(\rR)$ ranges over every integer in $[1,k]$ and $\kappa(\rW)$ from $O(1)$
to $\Theta(\min(k,m/w))$ (proved by explicit families in Supplementary Note~\ref{note:proofs-avoidance}). Writing
$\sigma(\rho,X)$ for predicted descriptor length at $X$'s density over actual, $\sigma$ is identically
one for \rB\ and \rS, at least $\max(d,\,1-d+\tfrac{1}{m})\ge\tfrac12$ for \rR, and unbounded above:
real rows can be arbitrarily cheaper than random rows of the same density and at most about twice as
expensive.

One impossibility follows, exactly. A selector reading density --- or cardinality, or any function of
them --- evaluates one fixed function on every row of that density, so by the range above it
mis-predicts $\kappa(\rR)$ by an unbounded factor on some rows at every density; a fixed representation
cannot exploit $\sigma$ either, having committed before $\sigma$ is knowable. That rules out a family
of selectors without choosing among the survivors, and two further readings do not follow: whether a
mis-predicting model costs more than selecting nothing depends on which cell it picks and that cell's
actual cost, which is measured (Fig.~\ref{fig:selection-cost}b), not derived; and a run or occupancy
count is $O(1)$ metadata that sees $\sigma$ perfectly and could feed a model. What is proved is that
selection must be gated on a statistic that sees structure rather than on cardinality --- structural
summary or direct kernel measurement is resolved by the measured selector comparison in
Fig.~\ref{fig:selection-cost}b.

\paragraph*{The null model is a floor, and it is the same floor three times.}
Independent placement at fixed density is the \emph{extremal} arrangement for every mechanism used
here, always in the same direction (proofs in Supplementary Note~\ref{note:proofs-avoidance}). For a representation,
Eq.~\eqref{eq:expected} sits within a factor of two of the worst achievable descriptor. For a zone map
of bin width $W_z$, scattering maximises the occupied-bin fraction, so the summary's null cost is an
upper bound on its cost and a lower bound on its value. For the prefix filter of the narrowed-contract
mode, the entire mechanism depends on the operands through one dimensionless quantity, and any ordering
that concentrates rare elements into prefixes makes real prefixes collide less often than scattered
ones of the same length. In each case the null is what the mechanism is worth on data with no
structure to find, real data deviates from it only in the mechanism's favour, and that deviation is
exactly the quantity the mechanism converts into avoided work.

\paragraph*{The mid-density regime selects bitmap throughput.}
At $s=\tfrac12$ the interval has width $m/2$ and excludes nothing, the cheapest list-based cell spends
$m/2$ probes against $\lceil m/w\rceil$ words, and Eq.~\eqref{eq:expected} offers neither \rR\ nor \rW\
an advantage: Eq.~\eqref{eq:headroom} falls below one. In the generic case the mid-density band belongs
to $\cell{\rB}{\rB}$, and a selector choosing it there is correct rather than defeated. This bounds what
cardinality determines absent structure, not what any kernel must spend --- call it a \emph{metadata}
or \emph{identifiability floor}, never an information-theoretic one --- and a structured row at exactly
this density is a counterexample to the unqualified form (Supplementary Note~\ref{note:proofs-avoidance}).
Zone-map occupancy remains informative in this band because it records arrangement that cardinality
does not.

\paragraph*{Why the two mechanisms are selected between rather than stacked.}
Composing a zone map with a representation is not the product of the two, because filtering
\emph{concentrates} what it passes on: conditioned on a bin of width $W_z$ being occupied its expected
local density is $d/p$, $p=1-(1-d)^{W_z}$, not $d$ --- no summary hands a kernel anything sparser than
one bit per bin (Supplementary Note~\ref{note:proofs-avoidance}). Charging the summary pass $c/W_z$ of a bitmap scan and taking
the operands independent, the composed cost relative to a plain bitmap is
\begin{equation}
\frac{c}{W_z} \;+\; p_A\,p_B\;\kappa\!\Big(\frac{d}{p}\Big),
\label{eq:compose}
\end{equation}
$\kappa\le1$ being the envelope of Eq.~\eqref{eq:expected} normalised against the bitmap. The
composition never costs more than the summary alone, yet is floored at $c/W_z$, below which a sparse
representation is cheaper outright --- three regimes, the middle one exactly $x e^{-x} > c/w$ in
$x = W_z d$ bits per bin, opening where the summary undercuts enumeration and closing where bins stop
being empty. This is the analytic form of the gating result of Fig.~\ref{fig:selection-cost}c, and it
models \emph{work}: it bounds where a filter has work to remove and says nothing about what one costs
when it removes nothing.

\paragraph*{Each mechanism's liability is two-sided, for unrelated reasons at the two ends.}
A summary is indexed by the universe while every compressed representation is indexed by the data, so
their costs must cross. Charging the summary $c/W_z$ of a bitmap scan as in Eq.~\eqref{eq:compose},
and a one-word-per-element representation $w\,d$, the low crossing is
\begin{equation}
d_{\mathrm{low}} \;=\; \frac{c}{w\,W_z} \;=\; 3.1\times10^{-5} \qquad (c=1,\; w=64,\; W_z=512),
\label{eq:dlow}
\end{equation}
below which reading the summary costs more than enumerating the operands outright (Supplementary Note~\ref{note:proofs-avoidance}). At the other end $p_Ap_B\to1$: no bin is empty, the summary proves nothing, and only its own
$c/W_z$ remains. A run- or fill-based representation crosses at its own density, wherever its
descriptor length falls below $m/W_z$, but that a crossing exists at all is generic. Prefix filtering
has the same shape in $y=p^2/m$: for $y\ll1$ the surviving fraction is $y$ itself, so pruning improves
quadratically as prefixes shorten, while for $y>1$ prefixes collide by default and the index earns
nothing. This is why no mechanism here is applied unconditionally, and why each gate is a two-sided
test rather than a threshold.

\paragraph*{A narrowed contract, and what two integers then settle.}
Every statement above returns the cardinality itself. A caller may instead ask a narrower question ---
report only the pairs with $J(A,B)\ge t$, or only those with $|A\cap B|\ge\tau$, or restrict attention
to operands whose size falls in a band --- and the narrowing is itself exploitable, because
Eq.~\eqref{eq:interval} already says what sizes permit:
\begin{equation}
J(A,B) \;=\; \frac{|A\cap B|}{|A\cup B|} \;\le\; \frac{\min(a,b)}{\max(a,b)},
\label{eq:sizebound}
\end{equation}
so $J(A,B)\ge t$ requires $\min(a,b)\ge t\,\max(a,b)$, and $|A\cap B|\ge\tau$ requires
$\min(a,b)\ge\tau$ (proof in Supplementary Note~\ref{note:proofs-avoidance}). The hypotheses are that $a,b>0$ and $t\in(0,1]$;
nothing is assumed about arrangement, about $m$, or about how either operand is stored. Both conditions
are necessary and neither is sufficient, so the test discards pairs and never reports one --- the
survivors still go to an exact kernel, at the cost of one comparison between two integers already held
as metadata, no element of either set read. This is the Swamidass--Baldi pruning bound
\cite{swamidass2007bounds} together with the size (length) filter of
\cite{arasu2006exact}, in the two-sided form standardised by \cite{mann2016empirical} and applied
jointly with prefix filtering to Tanimoto over binary vectors by \cite{anastasiu2016tapnn}, restated here to place it against the mechanisms above
rather than as a result of this paper.
The operational narrowing and complete thresholded construction are given in Supplementary
Note~\ref{note:threshold}.

\paragraph*{What that removes, and why a wider corpus helps it.}
Modelling a corpus by drawing cardinalities independently with $\log a$ uniform on $[0,\ln K]$ --- the
density $\Pr(a)\propto 1/a$, exactly the $1/i$ cardinality spectrum this paper's benchmark protocol
mandates (``Corpora, provenance and loaders''), not an approximation to it ---
$\min(a,b)/\max(a,b)\ge t$ becomes $|\log a - \log b|\le\gamma$, $\gamma=\ln(1/t)$, and for two
independent uniforms on an interval of length $L=\ln K$,
\begin{equation}
\Pr\big[\text{pair can qualify}\big] \;=\; 1-\Big(1-\frac{\ln(1/t)}{\ln K}\Big)^{2},
\qquad \ln(1/t)\le\ln K,
\label{eq:survive}
\end{equation}
with the absolute form $\big(1-\ln\tau/\ln K\big)^{2}$ (proof and worked values in Supplementary Note~\ref{note:proofs-avoidance}, which agree with a $10^{5}$-sample simulation to three decimals). Expanding for small $\gamma$
gives $2\gamma\!\int\!\phi^{2}$ for any density $\phi$ of $\log a$, so the governing quantity is the
\emph{effective log-spread} $1/\!\int\!\phi^{2}$ of the size distribution, of which
Eq.~\eqref{eq:survive} is the log-uniform case. Two hypotheses carry real weight. Cardinalities are
integers, and $a=b$ passes at every $t$, so the continuous form is optimistic where sizes are small.
And the favourable scaling in $K$ --- Eq.~\eqref{eq:survive} decreases monotonically in $K$ --- is a
property of the log-uniform family and not of heavy tails generally: under the alternative reading in
which \emph{ranked} sizes scale as $1/i$, the surviving fraction is exactly $1-t$, independent of $K$,
since the effective log-spread saturates at two nats however far the sizes are spread. The favourable
scaling in $K$ is therefore claimed for the spectrum this paper benchmarks, not for every skewed
corpus.

\paragraph*{The size bound does not multiply with the within-pair mechanisms.}
Acting at a coarser granularity is not sufficient for savings to compose. The size test decides
whether a pair is formed at all, while representation choice and the zone map act inside a pair that
is formed, so the granularity argument of Eq.~\eqref{eq:compose} predicts a product --- and the
product is not achieved, because both mechanisms read the same statistic. The survivors of
Eq.~\eqref{eq:sizebound} are the pairs with $a\approx b$, and every within-pair mechanism is cheapest exactly
when $a$ and $b$ are \emph{dissimilar}, since its cost tracks $\min(a,b)$ against a fixed
$\lceil m/w\rceil$. Conditioning raises the mean cost of a survivor by
\begin{equation}
\frac{\mathbb{E}\big[\min(a,b)\;\big|\;\min/\max\ge t\big]}{\mathbb{E}\,\min(a,b)}
\;=\; \frac{L^{2}\big(K(1-t)-\gamma\big)}{\gamma\,(2L-\gamma)\,(K-L-1)}
\;\xrightarrow[\;t\to1\;]{}\; \frac{\ln K}{2},
\label{eq:shortfall}
\end{equation}
and the composed work exceeds the product of the two savings by that same factor; the derivation, the
numerical verification and the decorrelation control that isolates the shared statistic as the cause
are in Supplementary Note~\ref{note:proofs-avoidance}. The composition is still strongly worth taking. What is
refuted is only the product; the general rule is that mechanisms compose multiplicatively when they act
on different waste \emph{and} are driven by different operand statistics, and sharing a statistic costs
the composition Eq.~\eqref{eq:shortfall}'s factor even across granularities. Values here are analytic,
computed over the model above at $m=10^{6}$, and count work rather than time.

\paragraph*{Choosing a representation is a different problem from choosing a size, and only one of
them is separable.}
Everything above prices representations; a deployed library must \emph{choose} between them, and the
choice has to be stated as an optimisation before it can be called optimal. Partition the universe into
chunks of $M$ bits, each stored in one representation, $n=M/w$ words to a chunk and $e$ bits to a
stored array element. Given a corpus, a distribution $\pi$ over the chunk pairs an application
evaluates, and an operation $\omega$, the problem is to choose an assignment $\rho$ of representations
to chunks minimising $\sum_{(x,y)\sim\pi} C_\omega(\rho(x),\mu_x;\rho(y),\mu_y)$ subject to a byte
budget, where $\mu$ is the metadata of Eq.~\eqref{eq:expected} and $C_\omega$ counts items touched at
measured per-item constants. Storage is a \emph{separable} objective --- minimised one container at a
time --- whereas the query cost is quadratic in the assignment, since there is no such thing as the
cost of a container, only the cost of a pair (Supplementary Note~\ref{note:proofs-choice}). And the incumbent's three
container decisions --- array against bitset, array against run, bitset against run --- are one rule,
$\arg\min$ over serialized size, so its array-to-bitset threshold is $\tau_{\mathrm{stor}} = M/e$, a
property of the format carrying no machine constant and no property of the workload: the exact optimum
of a separable objective, and why it can be a compile-time constant. What follows is what changes
when the objective is not separable.

\paragraph*{The crossover is a formula, and the two constants in the field are its two limits.}
Write $c_w$ for a vectorised word \texttt{AND}-plus-popcount, $c_p$ for a membership probe into a
bitmap and $c_m$ for a merge step, and $\theta_{pw}=c_p/c_w$, $\theta_{mw}=c_m/c_w$. Promotion of a
chunk of cardinality $k$ from array to bitset is favourable against a bitmap partner exactly above
\begin{equation}
\tau_{\cap} \;=\; \frac{c_w}{c_p}\cdot\frac{M}{w}\;=\;\frac{n}{\theta_{pw}},
\qquad\text{whence}\qquad
\frac{\tau_{\mathrm{stor}}}{\tau_{\cap}} \;=\; \frac{w}{e}\,\theta_{pw},
\label{eq:tauratio}
\end{equation}
and against an array partner above $(\theta_{pw}/\theta_{mw}-1)\,k_Y$ for a partner of size $k_Y$
(proof in Supplementary Note~\ref{note:proofs-choice}). Pricing a byte at $\lambda\ge0$ and minimising query cost plus
$\lambda\,\mathrm{bytes}$ gives a threshold that runs from $\tau_{\cap}$ at $\lambda=0$ to
$\tau_{\mathrm{stor}}$ as $\lambda\to\infty$: every value between the two is optimal at some price of a
byte, and $\tau_{\mathrm{stor}}=M/e$ is the unique threshold at which promotion is never paid for in
bytes. The threshold also carries the pair count, so under any byte budget the optimum is not a
function of the container at all. The measured input is the symmetric crossover \nm{cardinality at
which an array--array merge stops beating a bitmap--bitmap popcount}, which fixes $\theta_{mw}$;
substituted into Eq.~\eqref{eq:tauratio} at $M=2^{16}$, $w=64$, $e=16$ together with \nm{measured ratio
of a bitmap probe to a merge step} it gives \nm{the ratio of the two thresholds}, recovering the
divergence between the two objectives from measurements of the constants and no free parameters. Both
crossovers share $c_w$, so measuring them together over-determines the model and the probe-to-merge
ratio becomes a checkable prediction rather than an input.

\paragraph*{A fixed threshold is optimal, under five hypotheses, and each of them fails.}
The loose form is false, so state the hypotheses first: if the candidates are
only \rS\ and \rB, the operation is a count-only intersection, the per-item constants do not depend on
the resident footprint, there is no byte budget, and every chunk faces the same partner mix, then the
optimal assignment \emph{is} the fixed threshold $\tau_{\cap}$ of Eq.~\eqref{eq:tauratio} --- a machine
constant, independent of the corpus. Dropping any one hypothesis breaks it, with an exact witness for
each (Supplementary Note~\ref{note:proofs-choice}): admitting \rR\ reintroduces the impossibility above, specialised to the
container decision, since $\cell{\rR}{\rB}$ costs a strictly increasing function of run count while
$\cell{\rS}{\rB}$ does not depend on it; letting the partner mix vary makes the optimal threshold
\begin{equation}
\tau^{*}(\beta,\bar k)\;=\;\frac{\beta\,c_w n+(1-\beta)(c_p-c_m)\,\bar k}
                                 {\beta\,c_p+(1-\beta)\,c_m},
\label{eq:taustar}
\end{equation}
where $\beta$ is the pair-weighted fraction of a chunk's partners that are bitmaps and $\bar k$ the
mean cardinality of the array partners: a statistic of the workload rather than of the machine, and
self-referential --- an optimal fixed
threshold is a fixed point $\tau=\tau^{*}(\beta(\tau),\bar k(\tau))$, which exists but need not be
unique; admitting a byte budget makes the threshold its dual variable, sweeping the whole interval
above; changing the operation flips the matrix's sign, since cardinality queries have no operation
dimension at all (union, difference and symmetric difference are inclusion--exclusion over the
intersection) but a \emph{materialising} union with a bitmap operand costs $n$ words whatever the other
side is, so no single stored assignment serves a mixed workload. The fifth hypothesis --- that the
per-item constants do not depend on footprint --- is the one that is not analytic and the one the
measurements say dominates: $c_w$ is not a constant, footprint re-enters through it, and a work model
that prices the choice cannot price its consequence, the same divergence recorded under
Eq.~\eqref{eq:compose} and treated as a measurement question, not an analytic one.

\paragraph*{Metadata settles the decision wherever the decision matters.}
Two identifiability questions must not be merged. Eq.~\eqref{eq:width} says what cardinality
determines about the \emph{answer}, pinned only at the density extremes. What cardinality determines
about the \emph{decision} is different and more favourable: every cost in this paragraph and the two
before it is a closed form in $(k,r,\ell+g)$ and the constants, so a selector holding those three
integers per operand evaluates the ordering exactly, in $O(1)$, touching no operand data --- which is
what ``selection must be near-free'' means, and a consequence of the cost model rather than an
engineering aspiration (proof in Supplementary Note~\ref{note:proofs-choice}). Holding only $k$, the ordering is not
determined once \rR\ is a candidate; one extra integer restores it and more cardinality precision does
not. Suppose the constants are known within a factor $\eta\ge1$ --- the uncertainty of the modelled
costs, distinct from the $\gamma$ of the size bound above --- and let $\Lambda$ be the ratio of the two
cheapest modelled costs: if $\Lambda>\eta$ the ordering is determined; if $\Lambda\le\eta$ it is
not, and choosing the second-best costs at most $\eta$. So for every pair, either the metadata
settles the decision or the decision does not matter to within the precision of the constants, and the
undetermined set is an $\eta$-neighbourhood of the boundaries rather than a region of unbounded error.
The caveat that makes $\eta$ interesting: it is the uncertainty of the \emph{model}, and where the
model omits the memory system it is not small --- the derived reason to measure a candidate rather than
only to model it.

\paragraph*{Why an advantage survives repairing the kernel beneath it.}
A pairwise batch costs $C_{\mathrm{meta}}$ plus the cell cost summed over the pairs a summary does not
discard and the regions it does not exclude; those sets are functions of cardinality, bin occupancy and
the data, and of nothing else, so no change of representation moves either set (Supplementary Note~\ref{note:proofs-choice}).
Consequently, if a better kernel divides every cell cost by $\alpha$, the advantage over the same batch
without avoidance is invariant in $\alpha$ when the summary is free and erodes only through the
summary's share of the surviving cost: improving a kernel cannot subsume avoidance, because avoidance
changes which work is asked for and a kernel changes only what asked-for work costs. Which savings
multiply is settled by the same rule as Eq.~\eqref{eq:shortfall}: a zone map reads bin occupancy while
representation choice reads cardinality and run count, so those multiply; the size bound reads
cardinality, which representation choice also reads, so those do not, and a container-threshold gain
must not be multiplied into a size-bound gain. Quantitatively, against an incumbent already holding the
compute-optimal threshold of Eq.~\eqref{eq:tauratio}, the modelled advantage of adding a zone map of bin
width $W_z$ peaks at
\begin{equation}
G_{\max}\;=\;\frac{\theta_{pw}\,w}{2\sqrt{c\,W_z}}
\qquad\text{at}\qquad d^{*}=\frac{\sqrt{c}}{W_z^{3/2}},
\label{eq:gmax}
\end{equation}
valid on $(\theta_{pw}w\sqrt{c})^{2/3}\le W_z\le\theta_{pw}w$ and saturating at
$\sqrt{\theta_{pw}w/c}\,/2$ above it, and falling below one outside a window whose lower edge is
Eq.~\eqref{eq:dlow} with the per-item constants carried, $c/(\theta_{pw}wW_z)$ (proof in Supplementary Note~\ref{note:proofs-choice}). Eq.~\eqref{eq:gmax} is a \emph{null-model} ceiling on structureless data, and by the
one-sidedness of $\sigma$ above, real corpora can only fall below the null, so a measured advantage
exceeding it is evidence of structure and must be reported as such rather than as agreement. And
Eq.~\eqref{eq:gmax} excludes the narrowed contract entirely: Eq.~\eqref{eq:survive} supplies its own
factor, and by the rule above the two are not multiplied.

\paragraph*{The irreducible regime selects throughput.}
The mid-band above is not only where these results stop applying; it is where a different optimization
applies, and the two are complementary rather than competing. Where Eq.~\eqref{eq:headroom} falls below
one and $p_Ap_B\to1$, the work $\cell{\rB}{\rB}$ performs is, in the generic case, work that is going to
be performed, and what remains to optimize is the per-cycle throughput of its $\lceil m/(wV)\rceil$
vector \texttt{AND}-plus-popcount operations --- a quantity none of the results above speaks to, and one
the popcount literature has already driven close to the hardware limit \cite{mula2018popcount}.
Eq.~\eqref{eq:headroom} and Eq.~\eqref{eq:compose} locate the boundary between the two domains; they do
not rank them.


\subsection*{Representations and the per-cell cost model}

Each row --- one set $X_i$ over a shared universe of $m$ bits --- is stored in one of several
representations, each with a distinct asymptotic cost for computing an intersection
\emph{cardinality}, $|X_i \cap X_j|$, rather than the intersection itself. We define each
representation below, together with the cost of the cell it forms when paired with itself or with
another representation; every cost-model symbol used elsewhere in this paper is defined here and only
here.

Four representations form the symmetric pairing matrix. A dense bitmap (\rB) stores $m$ bits directly,
one per universe position, and costs $\Theta(m)$ to intersect against another dense bitmap by AND and
popcount, independent of either side's cardinality. A sorted array (\rS) stores the set's
$|X|$ set positions as sorted integers and costs $\Theta(|X|)$ against a bitmap, since each element is
resolved by a single random probe. A run container (\rR) stores the set as $r$ maximal runs
$(a,b)$ of contiguous set bits and, against a bitmap fitted with a rank index (below), costs
$\Theta(r)$ --- independent of run length. A fill-compressed word stream (\rW, an EWAH-style
encoding of literal words interleaved with run-length-coded fill words) costs
$\Theta(\ell + g)$ against a bitmap, where $\ell$ is the literal-word count and $g$ the fill-group
count, and its fills are skipped in constant time using the same rank index used for \rR. A fifth
view, the complement (\rC), does not add a representation to the symmetric matrix; it is a strategy,
described in the next subsection, for reusing whichever of \rB, \rS, \rR or \rW is cheapest on a
row's complement when that row sits above density one-half.

The cost of each symmetric cell is, in the notation above: $\cell{\rB}{\rB}=\Theta(m)$ (naive; made
sub-linear by the zone-map filter below); $\cell{\rB}{\rS}=\Theta(|S|)$; $\cell{\rB}{\rR}=\Theta(r)$
with a rank index; $\cell{\rB}{\rW}=\Theta(\ell+g)$; $\cell{\rS}{\rS}=\Theta(|A|+|B|)$ by merge or
$\Theta(|A|\log(|B|/|A|))$ by galloping search when $|A|\ll|B|$; $\cell{\rS}{\rR}=\Theta(|S|+r)$ by
merge or $\Theta(|S|\log r)$ by binary search into runs; $\cell{\rS}{\rW}=\Theta(|S|+\ell+g)$;
$\cell{\rR}{\rR}=\Theta(r_A+r_B)$ by run merge; $\cell{\rR}{\rW}=\Theta(r_A+(\ell+g)_B)$; and
$\cell{\rW}{\rW}=\Theta((\ell+g)_A+(\ell+g)_B)$. Every cell except $\cell{\rB}{\rB}$ is sub-linear in
the universe size $m$. The formal statement of what these costs imply --- that the null model is a
floor rather than a prediction, and the exact conditions under which no representation beats a bitmap
--- is given with full proofs in Supplementary Note~\ref{note:proofs-avoidance}.

\subsection*{The pairing matrix and its cells}

The paper reports the ten cells of the
symmetric representation matrix formed from \{\rB, \rS, \rR, \rW\}: \cell{\rB}{\rB}, \cell{\rB}{\rS},
\cell{\rB}{\rR}, \cell{\rB}{\rW}, \cell{\rS}{\rS}, \cell{\rS}{\rR}, \cell{\rS}{\rW}, \cell{\rR}{\rR},
\cell{\rR}{\rW}, \cell{\rW}{\rW}. The complement strategy, \cell{\rC}{\rB}, is an eleventh, asymmetric
strategy rather than a twelfth symmetric cell: when a row's density places it above one-half, its
complement is computed and intersected using whichever of the ten cells above is cheapest for that
complement, and the requested cardinality is recovered by inclusion--exclusion,
$|A \cap B| = |B| - |A^{c} \cap B|$, where $A^{c}$ is $A$'s complement against the shared universe.
This convention is used consistently throughout: any count of ``ten'' cells in this paper refers to
the symmetric matrix over \{\rB, \rS, \rR, \rW\}, and \cell{\rC}{\rB} is reported separately as a
strategy rather than folded into that count. Roaring is never constructed as a cell of this matrix; it
is used throughout as an external baseline only (below).

Each cell is implemented independently against a shared representation layer providing one view per
format (bitmap, sorted list, run, EWAH-fill and complement). Several of the asymmetric cells have more
than one implemented strategy --- for example \cell{\rB}{\rR} has a rank-indexed variant and a scalar
per-run-scan variant --- and each variant was raced against the independent oracle described below
before being admitted to any comparison. The variant reported for a given cell in Results is the one
selected as fastest for the density regime or corpus shape being reported, not a single fixed
implementation chosen in advance; the full record of implemented variants, including those that did
not win any regime, is not itemized in this paper (Code availability).

\subsection*{Zone-map and rank-index construction}

Two mechanisms let a kernel decide how much of a row it needs to visit, rather than visiting all of it
faster. To decide how much of a dense row a pairing needs to scan, we used a \emph{zone map}, a
bitmap of coarser granularity in which bit $k$ records whether bin $k$ of the underlying row --- a
contiguous block of \nm{zone-map bin width in bits} bits --- contains any set bit. To answer a run's
or a fill's contribution to a bitmap without touching the bits it spans, we used a \emph{rank
(prefix-popcount) index}, following the rank/select lineage of Jacobson, Clark and Vigna's
\texttt{rank9} \cite{jacobson1989succinct,clark1996compact,vigna2008broadword}.

The zone map is built by a single forward pass over the row, or is free at query time when the
underlying storage already carries block-occupancy metadata, as columnar formats commonly do; its
storage cost is \nm{zone-map storage cost as a fraction of the bitmap it summarizes} of the bitmap it
summarizes. For a pair, the bitwise AND of the two rows' zone maps, $\mathrm{occ}_A \wedge
\mathrm{occ}_B$, is itself a bitmap of $m/\nm{zone-map bin width in bits}$ bits, and its popcount is
the \emph{exact} number of bins the pairing needs to visit, not an estimate: a popcount of zero proves
the two rows disjoint without visiting either. The rank index is built by a single forward
prefix-popcount pass and stores \nm{rank-index storage cost as a fraction of the bitmap it indexes} of
the bitmap it indexes; for a run or fill spanning positions $[a,b)$ on the indexed side, its
contribution to the intersection is answered in constant time as
$\mathrm{rank}_B(b) - \mathrm{rank}_B(a)$, the number of set bits in $B$ up to $b$ minus the number up
to $a$ --- this formula is used, unchanged, by every cell that pairs a run- or fill-based
representation against an indexed bitmap.

Neither mechanism is applied unconditionally. The zone-map filter is consulted only ahead of cells
that would otherwise scan a dense representation in full (\cell{\rB}{\rB} and \cell{\rB}{\rW}); it is
not applied to \cell{\rB}{\rS} or \cell{\rB}{\rR}, which are already work-proportional to the sparse
side and for which a summary can remove nothing a direct probe does not already skip. Applied to a row
whose bins are uniformly occupied, the filter finds nothing to skip and its cost is pure overhead; the
policy that decides, per pair, whether the filter is worth consulting is described as a selection
policy in the next subsection, and its necessity is established as a result rather than assumed as a
procedure (Fig.~\ref{fig:selection-cost}c). This is a re-parameterization of established occupancy-summary mechanisms rather
than a new one: BitFunnel's hierarchical block index \cite{goodwin2017bitfunnel}, WAH/EWAH's own
uniform-fill words, and Lucene's sparse-block bitset all skip empty regions by a similar test; what is
specific to this setting is applying the same test as an \emph{exact} work-count for
intersection cardinality, where the summary's own popcount is the number of bins that must be visited
rather than a heuristic estimate of it. The cost model for this filter, the crossings that bound where
it pays, and the derivation of the bin width used here are given in
Supplementary Note~\ref{note:zonemap}.

\subsection*{Selection policies: per-pair, per-tile, and probe-and-commit}

Selection --- choosing which cell to run for a given pair --- is distinct from the zone-map and
rank-index mechanisms above, which decide how much of an already-chosen cell's work can be skipped;
the two are reported separately throughout. Three selection granularities were evaluated.

Under \emph{per-pair} selection, precomputed $O(1)$ metadata for each row --- cardinality, run count,
and per-row zone-map bin-occupancy count --- is consulted once for every pair $(i,j)$, and the cell
predicted cheapest under the cost model above is chosen. Under \emph{per-tile} selection, rows are
first partitioned into tiles of \nm{tile width in rows} rows by sorting or bucketing on density and run
count, so that rows sharing similar characteristics fall into the same tile; one cell choice is then
made per tile, from the same $O(1)$ metadata aggregated over the tile, and applied to every pair inside
it, hoisting the decision out of the innermost loop. Under \emph{probe-and-commit}, named here as the
$\varepsilon\!\to\!0$ special case of Micro Adaptivity in Vectorwise
\cite{raducanu2013microadaptivity}, a small number (\nm{count of probe pairs sampled per tile before
committing}) of a tile's pairs are timed under each candidate cell before any decision is made, and the
tile commits to whichever cell was empirically fastest on that sample for its remaining pairs. All
three policies are evaluated against an explicit runtime budget of \nm{stated selection-cost budget as
a percentage of total runtime, from the design target this paper adopts} of total runtime; a policy
that exceeds the budget is reported as failing it, not as a qualified pass.

Selection cost was measured by instrumenting the decision path itself separately from the kernel calls
it dispatches to, so that the reported cost of choosing is not conflated with the cost of computing.
\emph{Regret} against a bucket oracle is reported alongside selection cost for every online policy: the
oracle is the best cell for a given bucket of pairs chosen with full hindsight after every candidate
cell has been timed on it, and a policy's regret is its runtime on that bucket expressed relative to
the oracle's runtime on the same bucket. The same bucket oracle is used to report regret for an
always-\rB baseline and for a closed-form per-pair cost model evaluated from the metadata above but
without online timing, so that measured selection can be compared directly against both a
no-selection baseline and a modelled-selection baseline on the same metric.

Two distinct oracles are reported in this paper and are never interchanged. The \emph{bucket oracle}
defined above operates at bucket granularity and is the denominator for every regret figure. The
\emph{best-cell oracle} reported as an upper-bound column in Table~\ref{tab:corpora} is coarser: for a
whole corpus it takes the single cell whose measured end-to-end runtime on that corpus is lowest,
chosen after the fact and charged no selection cost. Because it may not vary its choice within a
corpus, it is not a lower bound on the bucket oracle's runtime; it is reported only as a bound on what
a per-corpus selector could achieve, and the deployed-versus-oracle gap in Table~\ref{tab:corpora} is
regret against that bound alone. The choice these policies approximate is itself an optimisation
problem with a characterisable optimum; it is stated and solved, with the conditions under which a
metadata-only decision is sufficient, in Supplementary Note~\ref{note:proofs-choice}.

\subsection*{Corpora, provenance and loaders}

\begin{sloppypar}
Three corpus families were used. The first is the set of corpora CRoaring's own benchmark harness runs
by default (the \texttt{real-roaring-datasets} collection), including \texttt{census1881},
\texttt{census-income}, \texttt{weather\_sept\_85}, \texttt{wikileaks-noquotes}, \texttt{uscensus2000}
and the \texttt{dimension\_0}\emph{xx} family, together with the row-sorted (\texttt{\_srt}) clustering
variant of each that the original Roaring papers also report. The second is a set of larger modern
graph corpora added beyond CRoaring's default set, drawn from public network-graph collections
(including \texttt{com-LiveJournal}, \texttt{com-Orkut} and \texttt{as-skitter}). The third is genomic:
two population-scale variant call sets, UShER SARS-CoV-2 and gnomAD chromosome 21, both stored
variant-major (one row per variant site, one bit per sample), and one haplotype-major encoding of 1000
Genomes Project phase 3 chromosome 20 (one row per haplotype, one bit per site) used specifically to
demonstrate that storage layout, not the underlying data, determines whether a corpus lands in the
sparse or the mid-density regime. Universe size and density are reported per corpus in
Table~\ref{tab:corpora} and are not duplicated in prose here; every corpus is drawn against the same competitively tuned
CRoaring baseline used throughout Results.

As a loader-correctness check, the \texttt{census1881} loader was validated by reproducing the
published corpus statistics reported for it (universe size and mean row cardinality) against the
values reported in the Roaring bitmap literature \cite{lemire2016roaring}, rather than trusting the
loader's own output uninspected.
\end{sloppypar}

\subsection*{Container census}

The container counts in Table~\ref{tab:census} are read from CRoaring's own introspection API,
\texttt{roaring\_bitmap\_statistics()}, and not inferred from our measurements. The census is taken
after \texttt{run\_optimize()} and \texttt{shrink\_to\_fit()}, so that it reports the containers the
library settles on once it has had the opportunity to choose run containers, matching the tuned
baseline every ratio in this paper is computed against. Four fields are recorded per corpus:
\texttt{n\_array\_containers}, \texttt{n\_bitset\_containers}, \texttt{n\_run\_containers}, and the
total container count. The threshold governing that choice is a vendor fact rather than a measurement
of ours: CRoaring stores a $2^{16}$-bit chunk as an array container unless the chunk's own cardinality
exceeds the compile-time constant \texttt{DEFAULT\_MAX\_SIZE}, whose value is
4096 (declared as
\texttt{enum \{ DEFAULT\_MAX\_SIZE = 4096 \}} at \texttt{roaring.h:2486}, and applied at insertion by
\texttt{array\_container\_try\_add(ac, val, DEFAULT\_MAX\_SIZE)} in \texttt{container\_add},
\texttt{roaring.h:5690}, in the pinned version below). The constant is not a tuning knob: it is an
\texttt{enum} with no override, it is enforced as a structural invariant by
\texttt{array\_container\_validate} and \texttt{bitset\_container\_validate}
(\texttt{roaring.c:7133}, \texttt{roaring.c:8263}), and the portable serialization format writes no
array-versus-bitset type tag at all, reconstructing the type from cardinality against this same
constant on read (\texttt{roaring.c:14421}, \texttt{roaring.c:14553}). The same fixed commitment prices a run container against a bitset by iterating every run and
charging for its length (\texttt{run\_bitset\_container\_intersection\_cardinality},
\texttt{roaring.c:10901}).

\subsection*{Benchmark protocol and cache-residency statement}

Every reported cell variant, including CRoaring, was compiled with the same compiler and flags,
allocated with the same alignment, and given the same warm-up treatment before timing, so that no
comparison advantages one implementation by construction. CRoaring was tuned in good faith before every
comparison: \texttt{run\_optimize()} and \texttt{shrink\_to\_fit()} were called before timing began,
matching what a competent user of the library would do, and every ratio against CRoaring reported in
this paper is computed against that tuned baseline, not against CRoaring's untuned default
construction.

Per-cell density and strategy sweeps (Figs.~\ref{fig:density-sweep} and~\ref{fig:strategy-decay}) raced every candidate variant over an identical,
pre-generated pair list in an identical order, so that the cache state each variant encountered at
call time was comparable across variants; each variant's differential correctness against the oracle
(below) was checked before any of its timings were accepted. A warm-up pass of \nm{warm-up iteration
count or criterion} preceded timed repeats. Wall-clock time was measured with a monotonic
per-run timer; where a hardware cycle count is reported, it is derived from wall time and a
frequency measured on the host at run time rather than assumed from a nominal clock speed, and is
labelled as a derived quantity rather than a hardware performance-counter reading. On Apple silicon
hosts, benchmark threads were requested onto the performance core cluster to avoid the heterogeneous
core migration that otherwise mixes two different microarchitectures' timings into one nominally
single-threaded measurement.

The working-set-size sweep behind the residency comparison (Fig.~\ref{fig:strategy-decay}b) varied corpus footprint from
\nm{smallest working-set size in the sweep} to \nm{largest working-set size in the sweep} per host,
and cache-residency boundaries (L1, L2 and shared/last-level cache) were identified from each host's
own reported cache geometry rather than assumed to be uniform across hosts (below). Every throughput
number reported anywhere in this paper is qualified by the cache-residency regime it was measured in,
named once here by regime (L1-resident, L2-resident, DRAM-resident) so that Results and Discussion can
use the regime name as shorthand rather than restate cache sizes at each mention.

\subsection*{Statistics and aggregation}

Two different aggregation conventions were used, for two different kinds of comparison, and neither is
substituted for the other anywhere in this paper. For controlled synthetic sweeps in which every
variant races the identical pair list under identical conditions (the density sweep, the
asymmetric-cell strategy comparison, the working-set-size sweep), the reported cost is the
\emph{minimum} over \nm{repeat count for synthetic per-cell sweeps} repeats, on the standard
microbenchmarking rationale that the minimum estimates the underlying cost with scheduler and
system-noise inflation removed, while the mean instead reports how noisy the run happened to be. For
the real-corpus comparison (Table~\ref{tab:corpora}), where repeats necessarily interleave across corpora and share a
host with other system activity over a longer measurement window, the reported ratio is instead the
\emph{median} over \nm{repeat count for real-corpus comparisons} interleaved repeats, reported together
with the observed range; no per-corpus ratio in this paper is quoted to two significant figures without
that range being available alongside it.

Run-to-run variance was largest for the corpora and hosts named explicitly in Results and in the notes to
Table~\ref{tab:corpora}; any corpus whose repeats disagreed beyond \nm{stated dispersion tolerance used to flag a corpus
as unquotable} was excluded from headline ranges as unquotable, and named as excluded rather than
silently dropped. On \nm{count of corpora, of the total corpus set, where the winning cell itself
changed between repeats} corpora, the winning cell reported in Table~\ref{tab:corpora} itself changed
between repeats; this is disclosed here once so that Table~\ref{tab:corpora} can report a single,
clearly labelled cell per corpus
without re-explaining the variance underlying it at every mention.

\subsection*{Correctness oracle and differential testing}

Every cell variant is checked against an independent scalar oracle that computes intersection
cardinality by direct bit-by-bit or element-by-element counting, sharing no code path with any of the
variants under test, before that variant is admitted to any timed comparison; a variant that returns a
wrong cardinality is not reported as fast regardless of its measured speed. Correctness is tracked at
three levels: an ABI-level regression suite (\nm{test\_storm check count} checks against the oracle,
\nm{test\_storm failure count, expected "0"} failures) that is compiled and linked as C to serve as
the public C ABI's own regression test; a cell-level differential suite (\nm{test\_cells check count}
checks, \nm{test\_cells failure count, expected "0"} failures) that exercises every implemented cell
variant against the oracle across a grid of densities, structures and alignments; and a larger
differential run, on the order of
\nm{cross-architecture differential check count per host family}, executed once per host family during
the cross-architecture measurement campaign. The public C ABI surface was verified independently with
\texttt{nm}: \nm{count of unmangled STORM\_-prefixed exported symbols} unmangled \texttt{STORM\_}
symbols and \nm{count of defined mangled (C++) symbols, expected "0"} defined mangled symbols, confirming
that the public interface is C-linkable as declared.

The cell-level suite is driven by a fuzz grid over \{row count, word count, density, within-row
structure (uniform, clustered, run-based), alignment\}, following the same axes used for every kernel
change in this project. The regression suite's own ability to fail was itself verified rather than
assumed: each of a set of known defects fixed earlier in the project was reverted individually and the
suite was confirmed to fail before being re-fixed, establishing that a red result from the suite is
informative and not merely a suite that always passes.

\subsection*{Hardware and toolchain}

Measurements were taken on \nm{count of microarchitectures used across the whole paper, expected
"four"} microarchitectures spanning two independently hand-written SIMD kernel paths --- one using ARM
NEON, the other using x86-64 AVX-512 with the \texttt{VPOPCNTDQ} extension --- plus a portable scalar
fallback used where neither hand-written path applies. No architecture-specific SVE or SVE2 kernel
exists in this project's kernel sources at the time of writing: the Arm Neoverse-class hosts used here
run the generic NEON path, on hardware that is SVE2-capable but for which no SVE2-specific
intrinsic code was written, and no result in this paper should be read as evidence about an SVE2
kernel. The same applies to AVX2 and SSE4.2: build options exist to gate an older, separate code path
(described below) to those instruction sets, but no pairing-matrix cell was written against them.

Each host is identified by \nm{per-host microarchitecture name, core/cluster configuration, and nominal
clock speed} and by its cache hierarchy (\nm{per-host L1d, L2, and shared/last-level cache sizes}),
cross-referenced against the cache-residency regimes named above rather than restated per measurement.
All code was compiled with \nm{compiler name and version per host}, using \nm{compiler flags used, and
whether native-architecture dispatch (\texttt{-march=native} or equivalent) was enabled or disabled for
a given reported number}. The CRoaring baseline used throughout is pinned at version
\nm{CRoaring version pin, expected "v4.7.2" pending re-verification against current \texttt{master}},
vendored as a source amalgamation within the project repository, so that the exact baseline is
reproducible rather than referred to generically as ``CRoaring.''

\subsection*{Data and code availability}

The benchmark harness, kernel sources, and per-corpus loader scripts described above are available in
the project repository under the Apache License, Version 2.0. The pairing matrix, the selector
and the tile pipeline described in this paper are not, at the time of writing, exposed through the
project's public C ABI; the public ABI currently surfaces an earlier, separate code path not used for
any measurement reported here. This paper reports a systematic measurement of the pairing-matrix
mechanism as implemented in the project's benchmark and kernel sources, not a released, ABI-stable
tool, and this is stated here as a fact about the present state of the code rather than a property that
is apologized for. Public real-world corpora were obtained from their original sources and reproduced
locally with the loader scripts referenced above; raw per-corpus timing data, including the repeats and
dispersion referenced in Statistics and aggregation, is deposited as Supplementary material.
